\section{Introduction}\label{sec:introduction}

Cyber security is a field where practical experience clearly helps users refine their skills. Lectures, documentation can help explain theory, but learners also require safe environments where they can experiment and investigate systems in order to make mistakes and learn from them. \gls{ctf} and cyber ranges are the current standard for supporting this style of hands on learning, however, they often depend on fixed challenges. Once these are solved, the usefulness of them is decreased, with solutions often being shared or memorised. VulnForge explores whether artificial intelligence can help produce varied environments that users can supplement their cyber security training with.

VulnForge aims to be an educational platform that is able to generate and deploy short lived vulnerable Linux environments for CTF style exercises. The main of this project is to allow a learner to choose a vulnerability category and difficulty level, launch a challenge, connect to it through SSH, exploit the vulnerability, submit a flag, and receive guidance from a chatbot if necessary. This solves the reliance on static, manually authored challenges. Using a \gls{llm} to aid to the challenge generation helps solve the aforementioned issues of static challenges where value can quickly be diminished.

The intended users of VulnForge are cyber security students, self directed learners, and educators who want repeatable, practical exercises without requiring manually authoring for each individual challenge. For learners, the benefit is repeated exposure to realistic scenarios in a controlled setting, for educators, it reduces the effort needed in creating variations of lab tasks. Broadly, we want to investigate how generative AI can be used in a responsible manner to aid cyber security education. It is not a replacement for conventional learning, but rather a supplemental tool for varied training material, with contextual support.

The scope of the project focuses on a prototype of a controlled environment, it is not intended to be a public cyber range, an enterprise security platform, or a replacement to the existing CTF platforms. Rather, it supports individual challenge attempts, basic progress tracking, and guided help. Features such as large scale management of classes, team modes, browser based terminals and cloud deployment are out of scope. This narrower focus allows us to hone in on the core technical feasibility of whether generative AI is able to generate vulnerable environments that can be deployed, validated, used in a learning environment, and cleaned up in a safe way.

The approach combines a web application, infrastructure automation, container orchestration, and a constrained use of generative AI. The web dashboard provides the user a way to authenticate, select a challenge, alter account settings, view challenge history, submit a flag, and interact with an educational chatbot. The system utilises a structured challenge specification, and Ansible to provision an Ubuntu 24.04 \gls{lts} \gls{lxc} container, through Proxmox. The vulnerable service is injected into this ephemeral environment, the learner is given SSH access to it, and validation checks are run to ensure the challenge is adequately prepared.

A key assumption behind this project is that generative AI cannot be trusted entirely, it requires guardrails and supervision. The model may produce invalid, unsafe, or unsuitable content, and we must treat this generation as so until it has been internally audited and validated. This means the validation schema, deployment restrictions, privilege separation and cleanup mechanisms are essential parts of the wider pipeline. This demonstrates that VulnForge is not just about whether an LLM can generate vulnerable code for a user to exploit, but whether said generation can be integrated into a safe, educational workflow.

The main outcome of this project would be a working prototype that is able to demonstrate the full lifecycle of an AI assisted cyber security training platform. It should generate a challenge, deploy an isolated environment, allow a learner to interact with it, support it through a chatbot tutor, validate submission, store challenge history and clean up resources afterwards. VulnForge attempts to show that this is technically feasible, but also highlights the limitations of this solution, and where careful decisions and validation is required. It turns challenge creation from a manually authored task to an automated pipeline that be constrained, tested, and evaluated. This dynamic generation can improve variety, but introduces questions around reliability, safety, exploitation and pedagogical quality. Therefore, VulnForge aims to demonstrate both the promise and risks of using generative AI to create practical cyber security educational environments.
